


\section{Scope and attribution}\label{scope-and-attribution}

This paper supplies a reusable interface for triangle constructions on
finite linear spaces. The construction is due to Delaby, Leemans and
Tranchida {[}DLT, Construction 1{]}. The common-neighbor reconstruction,
correlation extension and deleted-pencil description below reuse the
project's earlier \emph{Incidence compression}, §4. We give
complete proofs so that the present paper is an actual dependency of an
application, rather than a collection of references to an old combined
manuscript. The point-fiber synchronization in Theorem 2 also covers
two-point lines. The determinant criterion in Theorem 5 makes explicit
the semilinear alternative omitted from DLT's affine-plane
classification. Neither construction nor classical linear algebra is
claimed as a new subject. The accompanying application identifies
precisely which source statements these arguments recover or correct.

All sets in the geometric part are finite. A linear space S has a point
set P, a line set L, at least two points on each line, at least two
lines through each point, and a unique line through any two distinct
points. A flag is an incident pair (p,L). A uniform space has k points
on each line and r lines through each point; consequently v−1=r(k−1),
where v=\textbar P\textbar. We use ``automorphism'' for a
type-preserving incidence automorphism and ``correlation'' when
interchange of points and lines is allowed. A duality need not have
order two; a polarity does.

\section{1. A directed relation on
flags}\label{a-directed-relation-on-flags}

Write F for the flag set. Define the relation D by

\[
 (p,L)D(q,M)\quad\Longleftrightarrow\quad p\in M\text{ and }q\notin L.
\tag{1}
\]

Because q belongs to M, the conditions imply p≠q and M=pq≠L. If A and B
are the zero-one matrices for equality of point and equality of line on
F, then

\[
 D=(A-I)(B-I).
\tag{2}
\]

Indeed the only possible intermediate flag is (p,M). This observation
also proves that the integer product has entries zero or one.

The triangle incidence system Δ(S) has elements (f,i), where f∈F and
i∈Z/3Z. Distinct types i,i+1 are incident when fDg, respectively; the
reverse incidence is determined by symmetry. Its adjacency matrix, in
cyclic type order, is

\[
 \begin{pmatrix}0&D&D^{\mathsf T}\\D^{\mathsf T}&0&D\\
 D&D^{\mathsf T}&0\end{pmatrix}.
\tag{3}
\]

The cyclic shift τ(f,i)=(f,i+1) is therefore a triality. If p0,p1,p2 are
an ordered noncollinear triple, its chamber has flag of type i equal to
(pi,p(i−1)pi), with subscripts read cyclically. Conversely, (1) forces
these three joining lines in any chamber and excludes collinearity. A
panel between (p,L,0) and (q,pq,1) has exactly one completion for each
a∈L\{p\}: the third flag is (a,qa,2). Thus every flag extends to a
chamber, and every panel has k−1 completions in the uniform case. In
particular Δ(S) is firm precisely when k≥3 and is thin precisely when
k=3.

For a general rank-two input, the existence of a chamber of the literal
triangle system already forces a six-cycle in the input incidence graph:
its three point labels and three line labels are distinct, consecutive
incidences come from the chamber, and the required nonincidences
prohibit repeated labels. Consequently, if that nonempty triangle system
is a geometry, the input has gonality at most three. This recovers DLT
Proposition 3.3; it does not assert that every input of such gonality
gives a geometry.

The outgoing degree of the flag (p,L) is v−\textbar L\textbar, since
each q outside L gives the unique successor (q,pq). Hence flag
transitivity of Δ(S) forces every line of S to have the same size k: its
automorphism group is transitive on each vertex type. The identity
v−1=r\_p(k−1), obtained by partitioning P\{p\} among the lines through
p, then forces a constant r\_p=r. This proves uniformity without
assuming that the automorphisms of Δ(S) have already been reconstructed.

\section{2. Reconstruction from common
neighbors}\label{reconstruction-from-common-neighbors}

\textbf{Lemma 1 (point and line signatures).} For distinct flags f,g in
a uniform linear space, the following alternatives hold:

{\def\LTcaptype{none} % do not increment counter
\begin{longtable}[]{@{}lll@{}}
\toprule\noalign{}
Relation between f and g & Common successors & Common predecessors \\
\midrule\noalign{}
\endhead
\bottomrule\noalign{}
\endlastfoot
same point & (r−2)(k−1) & not needed \\
different points & 0 or k−2 & not needed \\
same line & not needed & (k−2)(r−1) \\
different lines & not needed & 0 or r−2 \\
\end{longtable}
}

For two flags (p,L),(p,M) at the same point, a common successor uses a
line through p different from L and M, and then one of its k−1 other
points. For flags at different points p,q, its line must be pq. If
neither forbidden line is pq, its point may be any of the k−2 points on
pq other than p and q; otherwise there are no choices. Reversing the
reasoning proves the last two rows. Distinct lines either meet in one
point or do not meet, so the reverse count has no further possibilities.
This proves the lemma.

\textbf{Theorem 2 (type-preserving reconstruction).} If r≥3 and k≥2,
restriction to the underlying flags identifies Aut(Δ(S)) with Aut(S).

The value (r−2)(k−1) is strictly larger than k−2. Thus the maximum
off-diagonal common-successor relation on each type identifies precisely
which flags have the same point. Every type-preserving automorphism φ of
Δ therefore induces a permutation fi of P on type i. Write

\[
 \phi(p,L,i)=(f_i(p),M_i(p,L),i).
\]

The point labels occurring among the successors of (p,L,i) are exactly
P\textbackslash L. Applying φ gives

\[
 P\setminus M_i(p,L)=f_{i+1}(P\setminus L),
 \qquad M_i(p,L)=f_{i+1}(L).
\tag{4}
\]

Since the image is a flag, fi(p) belongs to every f(i+1)(L) with L
through p. At least two distinct lines pass through p and their
intersection is \{p\}. Intersecting these images shows fi(p)=f(i+1)(p).
All three point permutations are therefore a single permutation f, and
(4) shows that f maps every line to a line and that φ is its diagonal
action on flags. The inverse satisfies the same condition, so f is an
automorphism of S. Conversely every automorphism of S preserves (1),
proving the claimed equality. The same argument, with the target point
set written P′, applies to any isomorphism Δ(S)→Δ(S′) between uniform
spaces with the same parameters k,r: equation(4) becomes
P′\textbackslash M\_i=f\_(i+1)(P\textbackslash L), and the same
intersections synchronize the three point bijections. Thus every such
triangle isomorphism is induced by one unique isomorphism S→S′. Notice
that the proof never needs k−2\textgreater0; two-point lines are
included.

If k,r≥3, both exceptional common-neighbor values in Lemma 1 are strict
maxima. A correlation reversing the cyclic type order exchanges the
successor and predecessor signatures and therefore exchanges the two
flag partitions. On each type it therefore induces a bijection from
points to lines and from lines to points. A point fiber and a line fiber
meet in a flag exactly when their labels are incident, so these two
bijections form a duality δi of S. In particular the action on type 0 is
induced by one duality δ0. The globally defined lift of δ0 reverses the
cyclic type order, since it transposes D. Composing the given
correlation with its inverse lift gives a type-preserving automorphism
which is the identity on type 0. Theorem 2 forces that automorphism to
be the identity on every type. Thus all three actions were induced by
δ0. No extra pointwise kernel has been suppressed in this
reconstruction.

\textbf{Theorem 3 (the complete correlation extension).} For k,r≥3 there
is a natural isomorphism

\[
 \operatorname{Cor}(\Delta(S))\simeq C_3\rtimes_\epsilon
 \operatorname{Cor}(S),
\tag{5}
\]

where ε records whether a correlation preserves or interchanges the two
types of S, and ε=−1 acts on C3 by inversion. If S has no duality, (5)
is C3×Aut(S). If S has a duality, the different exact sequence with
kernel Aut(S) and quotient S3 splits if and only if S has a polarity.

For completeness, a correlation δ of S acts on a flag by δ(p,L)=(δp,δL)
when it preserves types and by (δL,δp) when it interchanges them. Define

\[
 \widetilde\delta(f,i)=(\delta f,\epsilon(\delta)i).
\]

This is a homomorphism into Cor(Δ), fixes type 0, and conjugates τ to τ
raised to ε(δ). The reconstruction just proved shows that τ and all
these lifts generate the full correlation group. The two displayed
factors have trivial intersection, proving (5). A polarity supplies an
involutory lift of a type transposition and, together with τ, a
complement S3. Conversely any complement has such an involutory lift,
whose induced duality of S has order two. This proves the splitting
criterion. A polarity is not required for (5) itself; that qualification
belongs only to the latter exact sequence.

The chamber correspondence and Theorem 2 give the useful equivalence

\[
 \Delta(S)\text{ is flag-transitive}\quad\Longleftrightarrow\quad
 \operatorname{Aut}(S)\text{ is transitive on ordered noncollinear triples}.
\tag{6}
\]

The same assertion holds for any specified subgroup of Aut(S), using its
induced diagonal action.

\section{3. Connectivity and the deleted-pencil
residue}\label{connectivity-and-the-deleted-pencil-residue}

For k,r≥3, Δ(S) is connected. In a fixed type, flags with the same point
have a common successor and flags with the same line have a common
predecessor, by Lemma 1. These moves connect all flags because the
incidence graph of a linear space is connected: two points have their
joining line. Thus each type lies in one component, and an edge between
successive types joins the three components.

\textbf{Theorem 4 (residue interface).} The residue at (p,L,0) is the
bipartite incidence structure with parts

\[
 \{Q:Q\notin L\},\qquad
 \{(a,M):a\in L\setminus\{p\},\ M\ne L,\ a\in M\},
\tag{7}
\]

and incidence Q∈M. For the first part, the corresponding original vertex
is (Q,pQ,1); the second part already gives the original vertex of type
2. The residue in the Van Nypelseer construction used in L092 has
precisely the same description with a allowed to range over all of L.
Formula (7) is obtained by deleting one pencil, not by assuming that
this deletion preserves connectivity.

To verify (7), incidence with (p,L,0) forces a type-1 vertex to be
(Q,pQ,1) for Q outside L, and a type-2 vertex to be (a,M,2) with
a∈L\{p\} and M≠L. Their mutual incidence requires Q∈M and a∉pQ. The
latter condition is automatic, since a≠p and pQ meets L only in p.~This
proves both directions and the claimed literal identification.

In projective or affine dimension at least three, each edge in (7) stays
inside the plane generated by L and Q. Different such planes have
disjoint outside-point sets and disjoint allowed-line sets. Thus the
residue is a disjoint union of its plane residues. Whenever those plane
residues are connected, this is its exact component decomposition.

\section{4. Determinants classify affine-plane frame
orbits}\label{determinants-classify-affine-plane-frame-orbits}

We write vectors as rows. Let K be a finite field and H a subgroup
between SL(d,K) and ΓL(d,K), with d≥2. A semilinear element (A,σ) acts
as v↦v\^{}σ A. For the affine group V⋊H, translating the first point of
an ordered noncollinear triple to zero identifies its orbit with the
H-orbit of an ordered independent pair of vectors.

\textbf{Theorem 5 (frame-orbit criterion).} If d≥3, H is transitive on
ordered independent pairs. If d=2, its orbits on ordered bases are in
bijection with the orbits on K* of its induced determinant-semilinear
transformations

\[
 t\longmapsto \det(A)\,\sigma(t).
\tag{8}
\]

In dimension two, the affine group is transitive on ordered noncollinear
triples exactly when the action (8) is transitive. For a linear
H≤GL(2,K) this is equivalent to H=GL(2,K). For a semilinear H it need
not imply containment of GL(2,K).

For d≥3, extend either ordered independent pair to a basis. A linear map
sending one complete basis to the other may have determinant c≠1.
Multiply one of the target basis vectors outside the given pair by c⁻¹.
The resulting map has determinant one and still sends the required pair
correctly.

For d=2, arrange an ordered basis as the two rows of an invertible
matrix B. If B and C have equal determinant, there is exactly one matrix
M=B⁻¹C with BM=C, and det M=1. Thus SL(2,K) acts regularly on every
determinant fiber. A semilinear transformation sends B to B\^{}σ A and
its determinant to (8). These facts prove both implications of the orbit
correspondence: an H-orbit has one orbit of determinant values;
conversely first move a determinant to the desired value and then use
the contained SL group within that fiber. The transformations (8) form a
group because they are induced by an actual action on matrices. For a
linear H, the determinant image is a subgroup of K* acting by
multiplication. Transitivity forces the full image K*; together with its
contained kernel SL(2,K), this gives H=GL(2,K).

The statements about equal-determinant fibers, semilinear determinant
transport, and orbit saturation are formalized in
\texttt{formal/DeterminantFibers.lean}. The geometric
classification, finite-field subgroup identifications and connectivity
proofs are written mathematics, not consequences already formalized in
that file.

\section{5. A classical projective-line normalizer
interface}\label{a-classical-projective-line-normalizer-interface}

We record an elementary proof of the normalizer used by the circle
branch of the application, so that it is not hidden in a
higher-dimensional collineation theorem. This is a classical result, not
a novelty claim.

\textbf{Theorem 6.} For the natural action on P¹(K), where K is any
finite field,

\[
 N_{\operatorname{Sym}(\mathbf P^1(K))}
   (\operatorname{PSL}(2,K))=\operatorname{P\Gamma L}(2,K).
\tag{9}
\]

Let f normalize PSL. Since PGL is sharply three-transitive and itself
normalizes PSL, composing f with a projectivity lets us suppose that f
fixes ∞,0,1. The stabilizer of ∞ in PSL consists of x↦ax+b with a a
nonzero square (all nonzero a in even characteristic). Its translation
subgroup U=\{x↦x+b\} is its unique Sylow p-subgroup, since U is normal
and the scalar factor has order prime to p.~Thus f normalizes U. Because
f fixes0, conjugation sends the translation by b to the translation by
f(b), giving

\[
 f(x+b)=f(x)+f(b).
\tag{10}
\]

The stabilizer of ∞ and0 in PSL consists of the square multipliers.
Since f also fixes1, conjugating the multiplier by a square s gives the
multiplier by f(s), and hence f(sx)=f(s)f(x) for every nonzero square s;
the assertion also holds for s=0. In even characteristic every field
element is a square, so f is multiplicative. In odd characteristic every
t∈K is a sum of two squares: the set of squares including0 and its
translate t minus that set both have (q+1)/2 elements, so they
intersect. Write t=s+u with s,u squares. Additivity gives

\[
 f(tx)=f(sx)+f(ux)=(f(s)+f(u))f(x)=f(t)f(x).
\]

Thus f is a field automorphism in all cases. Undoing the initial
projectivity proves containment in PΓL. Conversely projectivities and
field automorphisms normalize PSL by their matrix action, proving
equality. The proof includes q=2,3: the same translation and square
arguments apply, without assuming simplicity of PSL. ∎

\section{References and dependency
boundary}\label{references-and-dependency-boundary}

{[}DLT{]} R. Delaby, D. Leemans and P. Tranchida, \emph{Geometries with
trialities arising from linear spaces}, Algebraic Combinatorics 9
(2026), 379--402, DOI10.5802/alco.482. Construction (1) and the
geometric targets are theirs.

{[}IC{]} \emph{Incidence compression}, repository manuscript, §4. Its
flag-matrix, correlation and deleted-pencil arguments are reused above.
The application must charge this complete rewritten dependency; citing
it does not make its proof cost disappear.

The group classification of locally two-transitive linear spaces remains
an explicitly retained application input. The complete Hermitian norm-shell
proof is included in Part III below. A
source correction, an independently verified new question answer, and a
recovered principal contract have different meanings and receive
separate accounting.

